\documentclass[conference]{IEEEtran}

\usepackage{cite}
\usepackage{amsmath,amssymb,amsfonts}
\usepackage{algorithmic}
\usepackage{graphicx}
\usepackage{textcomp}
\usepackage{xcolor}
\usepackage{booktabs}
\usepackage{tagging}
\usepackage{url}
\usepackage{listings}

\lstset{
  basicstyle=\footnotesize\ttfamily,
  breaklines=true,
  frame=single,
  columns=fullflexible,
}

\usepackage{sc26repro}

% Comment out both lines before final submission.
%\usetag{explanation}
%\usetag{example}

\begin{document}

\twocolumn[%
{\begin{center}
\Huge
Appendix: Artifact Description/Artifact Evaluation
\end{center}}
]

%%%%%%%%%%%%%%%%%%%%%%%%%%%%%%%%%%%%%%%%%%%%%%%%%%%%%
%  AD Appendix
%%%%%%%%%%%%%%%%%%%%%%%%%%%%%%%%%%%%%%%%%%%%%%%%%%%%%

\appendixAD

\section{Overview of Contributions and Artifacts}

\subsection{Paper's Main Contributions}

\begin{description}
\item[$C_1$] Structure-exploiting gradient computation for INLA:
a custom differentiation rule for the full objective function covering factorization, solve, and selected inversion of BTA precision matrices, computing exact gradients while exploiting block sparsity. Fusing stages of the forward computation reduces stored intermediates, cutting memory by approximately $2\times$.
\item[$C_2$] Extension to multivariate models:
gradient computation for jointly modeling $k$ correlated variables by exploiting the factored structure of the multivariate precision matrix blocks, separating per-variable and cross-variable hyperparameter contributions (benchmarked up to $k{=}3$, $d{=}15$).
\item[$C_3$] Multi-GPU distributed AD:
distributing the backward pass across GPUs via temporal-domain decomposition with nested dissection reordering, synchronizing boundary intermediates via MPI allgather and staging intermediates on CPU memory, enabling AD on models up to 1.9M latent variables.
\item[$C_4$] Comprehensive evaluation on up to 128 GH200 GPUs:
a practitioner-oriented analysis disentangling algorithmic gains from framework effects, evaluating when FD can match AD via scaling, and characterizing where time and memory concentrate in the gradient pipeline.
\end{description}

\subsection{Computational Artifacts}

\begin{description}
\item[$A_1$] \url{https://github.com/affifboudaoud/DALIA} (BSD-3-Clause)
\item[$A_2$] \url{https://github.com/affifboudaoud/serinv} (BSD-3-Clause)
\item[$A_3$] Benchmark suite (DOI to be assigned upon acceptance)
\end{description}

\begin{center}
\begin{tabular}{rll}
\toprule
Artifact ID  &  Contributions &  Related \\
             &  Supported     &  Paper Elements \\
\midrule
$A_1$   &  $C_1$, $C_2$, $C_3$ & Tables 1--2 \\
        &                       & Figures 1--8\\
\midrule
$A_2$   &  $C_1$, $C_2$ & Tables 1--2 \\
        &                & Figures 4--8\\
\midrule
$A_3$   &  $C_1$--$C_4$ & Tables 1--2 \\
        &                & Figures 3--8\\
\bottomrule
\end{tabular}
\end{center}


%%%%%%%%%%%%%%%%%%%%%%%%%%%%%%%%%%%%%%%%%%%%%%%%%%%%%%%%
\section{Artifact Identification}
%%%%%%%%%%%%%%%%%%%%%%%%%%%%%%%%%%%%%%%%%%%%%%%%%%%%%%%%

% ---- A1: DALIA/ADELIA ----
\newartifact

\artrel

$A_1$ is the ADELIA framework (extending DALIA) implementing INLA-based
Bayesian spatio-temporal inference with JAX automatic differentiation.
It directly implements $C_1$ (custom backward pass for BTA systems),
$C_2$ (multivariate coregional extension), and $C_3$ (distributed AD
via temporal-domain decomposition and mpi4jax).

\artexp

ADELIA's AD gradients should match JAX's built-in reverse-mode AD
(AD-Loop) to within $1.2 \times 10^{-7}$ relative error on models
that fit in memory, and match finite-difference gradients to within
FD truncation error ($O(h^2)$, $h{=}10^{-3}$) on all models.
The distributed two-phase implementation should match the AD-Loop
reference on GST-C3 with $P{=}2$ to within $6.8 \times 10^{-6}$.
Per-gradient speedups of $2.4$--$50.9\times$ over finite differences
are expected across the ten benchmark models.
On the four production-scale multivariate models, ADELIA should
converge to valid stationary points where FD stalls.

\arttime

Artifact setup: $\sim$30~min (conda environment, pip installs,
mpi4jax compilation, data download).
Gradient validation: $\sim$60~min on 1 GPU.
Full experiment reproduction: $\sim$28 GPU-hours.
Figure generation from reference CSVs: $<$5~min.

\artin

\artinpart{Hardware}

CSCS Alps supercomputer with NVIDIA GH200 Grace Hopper Superchip
nodes. Each GPU is a Hopper-based H100 with 96\,GiB HBM3,
paired with a 72-core ARM Neoverse V2 CPU (Grace) with
480\,GiB LPDDR5X memory via NVLink-C2C.
Inter-GPU communication uses HPE Slingshot-11 interconnect
(200\,Gb/s per direction).
Experiments use 1--128 GPUs.

\artinpart{Software}

\begin{tabular}{@{}ll@{}}
\toprule
\textbf{Package} & \textbf{Version} \\
\midrule
Python & 3.11 \\
JAX / JAXlib & 0.8.1 \\
CuPy & 13.6.0 \\
NumPy & $\geq$ 1.23 \\
SciPy & $\geq$ 1.10 \\
mpi4jax & 0.5.2+ (built from source) \\
Cray MPICH & 8.1.x \\
CUDA & 12.x \\
Pydantic & $\geq$ 2.0 \\
\bottomrule
\end{tabular}

\smallskip\noindent
A conda \texttt{environment.yml} is provided.
\texttt{mpi4jax} must be compiled against the system MPI
(Cray MPICH with CPU-staged transfers, no CUDA-aware MPI or NCCL);
see \texttt{INSTALL.md}.

\artinpart{Datasets / Inputs}

Input data for 10 benchmark models ($\sim$1.5\,GB total) are
downloadable via \texttt{scripts/download\_data.sh}.
Models span univariate synthetic (GST-S/M/L), real-world temperature
(GST-T), multivariate synthetic (GST-C2/C3), and production-scale
multivariate (SA1, AP1, WA1, WA2) with up to 1.9M latent variables
and $d{=}15$ hyperparameters.
Scaling studies use temporal variants of WA1
($n_t \in \{2, 32, 64, 128, 256\}$, $n_s{=}1{,}247$)
and spatial variants of WA2
($n_s \in \{72, 282, 1{,}119\}$, $n_t{=}48$).

\artinpart{Installation and Deployment}

\begin{lstlisting}[language=bash]
conda env create -f environment.yml
conda activate adelia-env
pip install -e DALIA/ -e serinv/
# Build mpi4jax against system MPI
export MPI4JAX_BUILD_MPICC=$(which mpicc)
pip install mpi4jax
\end{lstlisting}

\noindent
Due to a cuSolver/CuPy conflict, JAX must be initialized
before importing DALIA. All experiment scripts handle this
automatically. See \texttt{INSTALL.md} for cluster-specific
module loads and \texttt{LD\_LIBRARY\_PATH} configuration.

\artcomp

The workflow consists of three tasks:

\noindent
$T_1$~\textbf{Gradient validation} validates correctness by comparing
AD gradients against finite-difference and AD-Loop references.
Run per-model scripts in \texttt{validation/<model>/validate.sbatch}
or all models via \texttt{validation/validate\_all.sbatch}
(1 GPU, $\sim$60~min).
Checks: forward-value relative error $< 10^{-6}$,
gradient cosine similarity $> 0.999$,
max component-wise relative error $< 10^{-2}$ vs.\ FD.

\noindent
$T_2$~\textbf{Performance experiments} measure per-gradient wall-clock
time, scaling behavior, framework decomposition, and energy consumption.
Each experiment directory (\texttt{experiments/tab2\_*}, \texttt{fig2\_*},
\texttt{fig4\_*}, \texttt{fig5\_*}, \texttt{fig6\_*}, \texttt{fig7\_*})
contains its own \texttt{run.sbatch} and \texttt{README.md} with
resource requirements and submission instructions.
Each experiment runs at least 10 timed gradient evaluations
after 2 warmup iterations and reports mean $\pm$ 95\%
confidence intervals.

\noindent
$T_3$~\textbf{Figure generation} produces paper figures from experiment
output CSVs using plotting scripts in \texttt{plotting/scripts/}.
Reference CSVs are provided in \texttt{plotting/data/} so that figures
can be regenerated without re-running $T_2$.

\noindent
Dependencies: $T_1$ is independent.
$T_2$ experiments are independent of each other.
$T_3$ depends on $T_2$ (or uses pre-computed reference CSVs).

\artout

$T_1$ outputs per-model validation reports confirming gradient
correctness. When AD-Loop fits in memory (GST-S, GST-C2, GST-C3,
GST-M), ADELIA matches it to $< 10^{-7}$.
On large models, AD--FD agreement is within $5 \times 10^{-4}$
to $6 \times 10^{-3}$, consistent with FD truncation error.

$T_2$ outputs CSV files with per-gradient timing, peak GPU memory,
energy consumption, and scaling metrics.

$T_3$ outputs PDF figures corresponding to the paper:
Table~2 (strategy comparison),
Figure~4 (wall-clock breakdown),
Figure~5 (scaling study),
Figure~6 (framework decomposition),
Figure~7 (resource/energy efficiency),
Figure~8 (performance analysis).


% ---- A2: Serinv ----
\newartifact

\artrel

$A_2$ is the Serinv structured sparse solver library providing
the BTA factorization and selected inversion algorithms central
to $C_1$ and $C_2$.
ADELIA ($A_1$) depends on Serinv for its CuPy-based forward-pass
computations and the FD baseline.

\artexp

Serinv's sequential and distributed BTA algorithms should produce
numerically identical results to dense reference solvers on small
problems. On large problems, the structured algorithms enable
$\mathcal{O}(n \cdot b^3)$ complexity vs.\ $\mathcal{O}(N^3)$
for dense solvers ($N = n \cdot b$).

\arttime

Installation: $<$5~min (included in $A_1$ setup).
Serinv's unit tests: $\sim$10~min.

\artin

\artinpart{Hardware}

Same as $A_1$.

\artinpart{Software}

Serinv requires Python~$\geq$~3.9, NumPy, SciPy, and optionally
CuPy (GPU) and mpi4py (distributed).
All dependencies are included in $A_1$'s conda environment.

\artinpart{Datasets / Inputs}

Serinv's test suite generates random BTA matrices internally.
No external data is needed.

\artinpart{Installation and Deployment}

Installed as part of $A_1$: \texttt{pip install -e serinv/}.

\artcomp

Run the test suite:

\begin{lstlisting}[language=bash]
cd serinv && pytest tests/
\end{lstlisting}

\noindent
For distributed tests:

\begin{lstlisting}[language=bash]
mpirun -n 4 pytest --with-mpi tests_wrappers/
\end{lstlisting}

\artout

All tests should pass. Sequential tests verify factorization and
selected inversion against dense reference implementations.
Distributed tests verify that partitioned results match sequential
execution.


% ---- A3: Benchmark suite ----
\newartifact

\artrel

$A_3$ is the benchmark suite that produces all quantitative results
in the paper ($C_4$), exercising $A_1$ and $A_2$ across 10 models
and up to 128 GPUs to validate $C_1$, $C_2$, and $C_3$.

\artexp

The benchmark suite should reproduce:
(1) Table~2 showing that only ADELIA scales to million-variable models
while generic AD strategies (AD-Dense, AD-Loop, AD-Ckpt) OOM;
(2) $2.4$--$50.9\times$ per-gradient speedups (Fig.~4);
(3) $8$--$25\times$ AD speedups across temporal and spatial scaling (Fig.~5);
(4) $5$--$8\times$ energy savings even when FD is scaled to 128 GPUs (Fig.~7).
Exact numerical values may vary with hardware and software versions,
but relative rankings and order-of-magnitude results should be
consistent.

\arttime

Table~2: $\sim$2~h (1 GPU).
Figure~4: $\sim$8~h (1--4 GPUs).
Figure~5: $\sim$6~h (1 GPU).
Figure~6: $\sim$3~h (1--4 GPUs).
Figure~7: $\sim$12~h (4--128 GPUs).
Figure~8: $\sim$1~h (4 GPUs).
Total: $\sim$28 GPU-hours.

\artin

\artinpart{Hardware}

Same as $A_1$. Figure~7 (resource efficiency) requires up to
128 GPUs for the FD scaling baseline.

\artinpart{Software}

Same as $A_1$, plus matplotlib for figure generation.

\artinpart{Datasets / Inputs}

Downloaded via \texttt{scripts/download\_data.sh} ($\sim$1.5\,GB).
See the benchmark model table in $A_1$.

\artinpart{Installation and Deployment}

Included in the artifact repository.
Experiment scripts source \texttt{experiments/common/env\_setup.sh}
for environment configuration.

\artcomp

Each experiment directory contains \texttt{run.sbatch},
\texttt{run\_inner.sh} (environment + srun wrapper), and either
a \texttt{run.py} or per-model subdirectories with their own scripts.
Each directory's \texttt{README.md} documents parameters and
resource requirements.

Plotting scripts in \texttt{plotting/scripts/} read CSV outputs
and produce PDF figures. Reference CSVs from the paper are
provided in \texttt{plotting/data/}.

\artout

Each experiment writes timing CSVs and console output.
Plotting scripts produce PDF figures matching the paper:
\begin{itemize}
  \item \texttt{plot\_convergence\_ap1.py} $\rightarrow$ Figure~3
  \item \texttt{plot\_structure\_comparison.py} $\rightarrow$ Table~2
  \item \texttt{plot\_pipeline\_wallclock\_merged.py} $\rightarrow$ Figure~4
  \item \texttt{plot\_scaling\_study.py} $\rightarrow$ Figure~5
  \item \texttt{plot\_framework\_decomposition.py} $\rightarrow$ Figure~6
  \item \texttt{plot\_resource\_energy\_efficiency.py} $\rightarrow$ Figure~7
  \item \texttt{plot\_performance\_analysis.py} $\rightarrow$ Figure~8
\end{itemize}

\end{document}
